\documentclass[10pt]{article}

\usepackage[utf8]{inputenc}

\usepackage[T1]{fontenc}

\usepackage{amsmath}

\usepackage{amsfonts}

\usepackage{amssymb}

\usepackage[version=4]{mhchem}

\usepackage{stmaryrd}



\begin{document}

Лum.: [1] С а в е л о в А. А., Плоские кривые, М., 1960.



IIIТРМА ТЕОРЕМА: если



Д. Д. Соколов.



\[

f_{0}(x), f_{1}(x), \ldots, t_{s}(x)

\]



\begin{itemize}

  \item ряд Шітурма для отрезка $[a, b], a<b$ и $w(x)$ - число перемен знака в ряде $(*)$ в точке $x \in[a, b]$ (значения, равные нулю, не учитываются), то число различных корней функции $f(x)$ на отрезке $[a, b]$ равно разности $w(a)-w(b)$.

\end{itemize}



Р я д о м III т у р м а наз. последовательность деїствительных функций $(*)$, вепрерывных на отрезке $[a, b]$ и имеющих на этом отрезке конечное число корней, и такая, что



\begin{enumerate}

  \item $f_{0}(a) f_{0}(b) \neq 0$



  \item $f_{s}(x) \neq 0$ на $[a, b]$



  \item из $f_{k}(c)=0$ для нек-рого $k(0<k<s)$ и данного $c$ из $[a, b]$ следует $f_{k-1}(c) f_{k+1}(c)<0$,



  \item из $f_{0}(c)=0$ для данного $c(a<c<b)$ следует, что для достаточно малого $\varepsilon>0$



\end{enumerate}



\[

\begin{array}{ll}

f_{0}(x) f_{1}(c)<0 & (c-\varepsilon<x<c) ; \\

f_{0}(x) f_{1}(c)>0 & (c<x<c+\varepsilon) .

\end{array}

\]



Теорема доказана III. Штурмом [1], к-рый указал также следующий метод построения ряда П̈турма для многочлена $f(x)$ с действительными коәффициентами, ве имеющего кратных корней: $f_{0}(x)=f(x), f_{1}(x)=f^{\prime}(x)$ и, если уже построены многочлены $f_{0}(x), \ldots, f_{k}(x)$, то за $f_{k+1}(x)$ надо брать с обратным знаком остаток при делении $f_{k-1}(x)$ на $f_{k}(x)$. При этом $f_{\mathcal{S}}(x)$ будет константой, отличной от нуля.



Jum.: [1] S t u r m J. Ch., «Bull. de Férussac», 1829, t. 11; [2] К у р о ші А. Г., Курс высшей алгебры, 11 изд., М., 1975. ШТУРМА - ЛИУВИЛЛЯ ЗАДАЧА - задача, порожденная на конечном или бесконечном интервале $(a$, $b)$ изменения переменной $x$ уравнением



\[

-\frac{d}{d x}\left(p(x) \frac{d y}{d x}\right)+l(x) y=\lambda r(x) y

\]



и нек-рыми граничными условиями, где $p(x)$ и $r(x)$ положительны, $l(x)$ действительна, а $\lambda-$ комплексный параметр. Начало глубокому изучению этой задачи положили III. Штурм (Ch. Sturm) и Ж. Лиувилль (J. Liouville). Понятия и методы, зародившиеся в процессе изучения Ш.- Л. з., сыграли больпую роль в развитии многих направлений математики и физики. Она была и остается постоянным источником новых идей и задач для спектральной теории операторов и смежных вопросов анализа. Особое значение приобрела она в последнее время после открытия связи с нек-рыми нелинейными әволюционными уравнениями математич. физики.



Если $p(x)$ дифференцируема, а $p(x) r(x)$ - дифференцируема дважды, то уравнение (1) с помощью подстановки сводится к виду (см. [1])



\[

-y^{\prime \prime}+q(x) y=\lambda y .

\]



Принято различать регулярные и сингулярные задачи. ШІ.- Л. 3. для уравнения (2) наз. р е г у л я рн 0 й, если интервал $(a, b)$ изменения переменной $x$ конечен и если функция $q(x)$ суммируема во всем интервале $(a, b)$. Если же интервал $(a, b)$ бесконечен или $q(x)$ несуммируема (или и то и другое), то эта задача наз. с и н $\mathbf{⿴}$ у л я р н о ї.



Ниже рассматриваются в отдельности следующие случаи: 1) интервал $(a, b)$ конечен, в этом случав, не нарушая общности, можно считать, что $a=0$ и $b=\pi$; 2) $a=0, b=\infty$; 3) $a=-\infty, b=\infty$.



\begin{enumerate}

  \item Рассматривается задача, порожденная на сегменте $[0, \pi]$ уравнением $(2)$ и разделенными граничными условиями

\end{enumerate}



\[

y^{\prime}(0)-h y(0)=0, y^{\prime}(\pi)+H y^{\prime}(\pi)=0,

\]



где $q(x)$ - действительная суммируемая на сегменте $[0, \pi]$ функция, $h$ и $H$ - произвольные конечные или бесконечные фиксированные действительные числа, $\lambda$ - комплексный параметр. Если $h=\infty \quad(H=\infty)$, то первое (второе) условие в (3) заменяется условием $y(0)=$ $=0 \quad(y(\pi)=0)$. Для определенности далее предполагается, что числа, участвующие в граничных условиях, конечны.



Число $\lambda_{0}$ наз. с о б с т в е в в ы м 3 н а ч е н и е м задачи (2), (3), если при $\lambda=\lambda_{0}$ уравнение (2) имеет нетривиальное решение $y_{0}(x) \not \equiv 0$, удовлетворяющее граничным условиям (3); при этом функция $y_{0}(x)$ наз. с о б с т в е н н о й ф у нк ц и е й, соответствующей собственному значению $\lambda_{0}$.



Собственные значения граничной задачи (2), (3) действительны; каждому собственному значению соответствует единственная линейно независимая собственная функция (в силу действительности $q(x)$ и чисел $h, H$ собственные функции задачи (2), (3) можно выбрать действительными); собственные функции $y_{1}(x)$ и $y_{2}(x)$, соответствующие различным собственным значениям, ортогональны, т. е. $\int_{0}^{\pi} y_{1}(x) y_{2}(x) d x=0$.



Существует неограниченно возрастающая последовательность собственных значений $\lambda_{0}, \lambda_{1}, \lambda_{2}, \ldots$ граничной задачи (2), (3); при этом собственная функция $y_{n}(x)$, соответствующая собственному значению $\lambda_{n}$, имеет ровно $n$ нулеї в интервале $(0, \pi)$.



Пусть $W_{2}^{m}[0, \pi]$ - пространство Соболева, состоящее из заданных на сегменте $[0, \pi]$ комплекснозначных функций, к-рые имеют $m-1$ абсолютно вепрерывных производных и производную порядка $m$, суммируемую на сегменте $[0, \pi]$. Если $q(x) \in W_{2}^{m}[0, \pi]$, то собственные значения $\lambda_{n}$ граничной задачи (2), (3) при больших $n$ удовлетворяют асимптотич. равенству (см. [4]):



\[

\sqrt{\lambda_{n}}=n+\sum_{1 \leqslant 2 j+1 \leqslant m+2} \frac{c_{2 j+1}}{n^{2 j+1}}+

\]



$+\frac{(-1)^{m-1}}{2^{m+2}}\left(S_{m}(n)+\frac{\tilde{S_{m n}}(n)}{n}\right) \frac{1}{n^{m+1}}+\frac{\delta_{n}}{n^{m+2}}+\frac{\varepsilon_{n}(h, H)}{n^{m+3}}$,



где $c_{2 j+1}$ - независимые от $n$ числа,



\[

c_{1}=\frac{1}{\pi}\left(h+H+\frac{1}{2} \int_{0}^{\pi} q(t) d t\right),

\][



S\_\{m\}(n)=\textbackslash frac\{2\}\{\textbackslash pi\} \textbackslash int\_\{0\}\^{}\{\textbackslash pi\} q\^{}\{(m)\}(t) \textbackslash sin \textbackslash left\{2 n t-\textbackslash frac\{\textbackslash pi\}\{2\}(m+1)\textbackslash right\} d t,

]



$\tilde{S}_{m}(n)=\frac{2}{\pi} \int_{0}^{\pi} q^{(m)}(t)\left(2 h-c_{1} t\right) \sin \left\{2 n t-\frac{\pi}{2}(m+2)\right\} d t$, $\delta_{n}$ не зависит от $h, H$ и



\[

\sum_{n=0}^{\infty}\left|\delta_{n}\right|^{2}<\infty, \sum_{n=0}^{\infty}\left|\varepsilon_{n}(h, H)\right|^{2}<\infty .

\]



Отсюда, в частности, следует, что если $q(x) \in W_{2}^{1}[0, \pi]$, то



\[

\lambda_{n}=n^{2}+c+\frac{\gamma_{n}}{n}

\]



где



\[

c=\frac{2}{\pi}\left(h+H+\frac{1}{2} \int_{0}^{\pi} q(t) d t\right), \sum_{n=0}^{\infty}\left|\gamma_{n}\right|^{2}<\infty .

\]



Поэтому ряд $\sum_{n=0}^{\infty}\left(\lambda_{n}-n^{2}-c\right)$ сходится. Его сумма наз. ре гуля ри зо в а вным с лед о м задачи (2), (3) (см. [13]):



\[

\sum_{n=0}^{\infty}\left(\lambda-n^{2}-c\right)=\frac{q(0)+q(\pi)}{4}-\frac{(h+H)^{2}}{2}+h H-\frac{c}{2} .

\]



Пусть $v_{0}(x), v_{1}(x), \ldots-$ ортонормированные собственные функции задачи (2), (3), соответствующие собственным значениям $\lambda_{0}, \lambda_{1}, \ldots$. Для каждой функции $f(x) \in$





\end{document}